\documentclass[10pt]{article}
\usepackage[margin=1.8cm]{geometry}
\usepackage{booktabs}
\usepackage{amsmath}
\usepackage[hidelinks]{hyperref}
\setlength{\parskip}{3pt}
\setlength{\parindent}{0pt}

\begin{document}

\begin{center}
{\Large Kangaroo arms and the BeyondMimic action scale}\\[3pt]
Unitree G1 vs PAL Kangaroo. Summary of the SuperMimic gain A/B and the 7-DoF WBC wrist report
(M.~Actis). pal\_mjlab, October 2026.
\end{center}

\section{The rule}

BeyondMimic, and SuperMimic after it, derive gains and action scale per joint from the motor alone:
\[
K_p = I_a\,\omega^2,\quad K_d = 2\zeta I_a\,\omega,\quad \omega = 2\pi f,\ \zeta = 2,\qquad
q_{sp} = q_0 + \alpha\,a,\quad \alpha = c\,\frac{\tau_{\max}}{K_p},\ c = 0.25 .
\]
$I_a$ is the armature (reflected motor inertia), $\tau_{\max}$ the effort limit, $q_0$ the default
pose, $a$ the network output and $\alpha$ the setpoint offset per action unit (rad/unit). mjlab sends
$q_{sp}$, minus the calibration offset, to a MuJoCo position actuator:
$\tau = \mathrm{clip}\big(K_p(q_{sp}-q) - K_d\dot q,\ \pm\tau_{\max}\big)$. The policy only moves the
setpoint; the one spring pulls the joint toward it.

\section{What an action unit means}

The same law in action units, using $q_{sp} = q_0 + \alpha a$:\quad
$\tau = \underbrace{c\,\tau_{\max}\,a}_{\text{action}} - \underbrace{K_p(q-q_0)}_{\text{where }q\text{ sits vs }q_0}
- \underbrace{K_d\,\dot q}_{\text{damping}}$.
This is bookkeeping, not a second spring: actions are measured from $q_0$, so a joint resting at $q$
(setpoint on the joint, $\tau \approx 0$) needs $a = (q-q_0)/\alpha$.
\begin{itemize}
  \setlength{\itemsep}{1pt}
  \item \textbf{Torque per unit:} $\partial\tau/\partial a = c\,\tau_{\max}$, i.e.\ 25\% of the effort
  limit, at a fixed state only.
  \item \textbf{Holding} away from $q_0$ is a large action with almost no torque, so $|a|>4$ is not
  saturation. WBC example, arm joint 1 at
  $a=4.5$: $+56$\,Nm action term $-\,49$\,Nm $(q-q_0)$ term $=7$\,Nm net.
  \item \textbf{Moving} takes $\Delta a \approx \dot q\,\Delta t/\alpha$ per step. \textbf{Exploring}
  with noise $\sigma$ moves the setpoint by $\sigma\alpha$.
\end{itemize}
So $K_p/\tau_{\max} = c/\alpha$ is the \textbf{exchange rate between joint motion and action units}.
Torques don't depend on it, but the action-rate penalty, the exploration noise and the output range
do. It follows from the hardware and the frequency:
\[
\frac{K_p}{\tau_{\max}} = \frac{I_a}{\tau_{\max}}\,\omega^2
\quad\Rightarrow\quad
\frac{(K_p/\tau_{\max})_{\text{Kangaroo}}}{(K_p/\tau_{\max})_{\text{G1}}}
= \frac{(I_a/\tau_{\max})_{\text{Kangaroo}}}{(I_a/\tau_{\max})_{\text{G1}}}\cdot
\Big(\frac{f_{\text{Kangaroo}}}{f_{\text{G1}}}\Big)^{2}.
\]

\section{Unitree G1 vs Kangaroo}

\begin{center}
\footnotesize
\setlength{\tabcolsep}{4pt}
\begin{tabular}{llrrrrrrr}
\toprule
Joints & Motor & $I_a$ & $\tau_{\max}$ & $f$ & $K_p$ & $I_a/\tau_{\max}$ & $K_p/\tau_{\max}$ & $\alpha$ \\
 & & (kg\,m$^2$) & (Nm) & (Hz) & (Nm/rad) & ($10^{-3}$\,s$^2$) & (1/rad) & (rad/unit) \\
\midrule
\multicolumn{9}{l}{\textit{Unitree G1, BeyondMimic (mjlab)}} \\
Shoulders, elbow, wrist roll & 5020 & 0.0036 & 25 & 10 & 14.3 & 0.14 & 0.57 & 0.44 \\
Hip pitch/yaw, waist yaw & 7520-14 & 0.0102 & 88 & 10 & 40.2 & 0.12 & 0.46 & 0.55 \\
Hip roll, knee & 7520-22 & 0.0251 & 139 & 10 & 99.1 & 0.18 & 0.71 & 0.35 \\
Waist pitch/roll, ankles & 2$\times$5020 & 0.0072 & 50 & 10 & 28.5 & 0.14 & 0.57 & 0.44 \\
Wrist pitch/yaw & 4010 & 0.0043 & 5 & 10 & 16.8 & 0.85 & 3.36 & 0.075 \\
\midrule
\multicolumn{9}{l}{\textit{PAL Kangaroo (PR \#126 constants)}} \\
Leg 1--3 (hip) & & 0.20--0.50 & 80--230 & 3.5 & 97--242 & 1.4--3.1 & 0.70--1.51 & 0.17--0.36 \\
Leg 4, 5 (ankle) & & 0.05--0.10 & 82--140 & 3.5 & 24--48 & 0.6--0.7 & 0.29--0.35 & 0.72--0.85 \\
Arm 1, 2 + waist & S+, gear 121 & 0.483 & 50 & 3.5 & 233.7 & 9.7 & \textbf{4.67} & \textbf{0.053} \\
Arm 3--5 & S$-$, gear 101 & 0.305 & 25 & 3.5 & 147.5 & 12.2 & \textbf{5.90} & \textbf{0.042} \\
Arm 6, 7 & XS, gear 101 & 0.203 & 25 & 3.5 & 98.2 & 8.1 & \textbf{3.93} & \textbf{0.064} \\
Arms + waist at 1\,Hz & as above & as above & as above & 1 & 8.0--19.1 & as above & 0.32--0.48 & 0.52--0.78 \\
\bottomrule
\end{tabular}
\end{center}

\textbf{Arm vs arm.} Kangaroo's geared arm motors have 56--134$\times$ G1's armature for only
1--2$\times$ its effort, so $I_a/\tau_{\max}$ is \textbf{56--85$\times$} larger. At 3.5\,Hz instead of
10\,Hz, $\omega^2$ shrinks by $8.2\times$, so $K_p/\tau_{\max}$ ends up \textbf{7--10$\times$} G1's:
0.57 on G1 against 3.9--5.9 on Kangaroo. Matching G1 needs $f \approx 10/\sqrt{56\text{--}85} \approx 1.1$--$1.3$\,Hz.
Kangaroo's legs are already close to G1 (0.29--1.51). The only G1 joints in Kangaroo's arm regime are the two
5\,Nm wrist motors (3.36), whose motion barely changes the tracked bodies. On Kangaroo the regime covers
14 arm joints plus the waist.

\section{What it does to learning}

\begin{itemize}
  \setlength{\itemsep}{1pt}
  \item \textbf{Wrist exploration (WBC report).} At $\alpha \approx 0.04$--$0.06$, exploration barely moves
  joints 5--7. Their effect, end-effector orientation, has a reward flat near 0, so PPO gets no signal and
  leaves the wrist at $q_0$.
  \item \textbf{Action size and action-rate penalty (SuperMimic).} The arm references need $|a|$ = 7.8
  (median) to 40.8 (p95). The penalty $-0.1\lVert\Delta a\rVert^2$ scales with $(K_p/\tau_{\max})^2$ for the
  same motion: 28\% of the maximum per-step reward at p95 at 3.5\,Hz, against 0.2\% at 1\,Hz.
  \item \textbf{Not the plant.} Torque saturation is near zero in both studies.
\end{itemize}

\section{Evidence}

\begin{center}
\small
\begin{tabular}{p{3.4cm}p{4.6cm}p{8.6cm}}
\toprule
Study & Change & Result \\
\midrule
WBC, 7-DoF (M.~Actis) & 3.5\,Hz gains kept; $\alpha$ on joints 5--7 only, $\times4$ ($c{=}1$) and
$\times16$ ($c{=}4$) &
Orientation reward rises close to the 1\,Hz run; position and root unchanged. The scale alone is the
main factor for the wrist. \\
SuperMimic PPO-Sim, 1296 clips & Arms + waist at 1\,Hz vs 3.5\,Hz (gains and $\alpha$ together),
1000 iterations &
Arm error 21.2$\to$15.3\,cm (better on 90\% of clips), legs 11.3$\to$8.6\,cm (95\%), tracking return
0.53$\to$0.62, saturation 0.8\%$\to$0.04\%. At 3.5\,Hz the arm error grew during training. \\
\bottomrule
\end{tabular}
\end{center}
1\,Hz multiplies $\alpha$ by $12.25$, the same as $c \approx 3.06$ at 3.5\,Hz, which lies within the WBC
report's recommended $c = 3$--$4$ for the wrist.

\section{Options and rule of thumb}

\begin{enumerate}
  \setlength{\itemsep}{1pt}
  \item \textbf{Arms and waist at $\sim$1\,Hz} (current SuperMimic baseline). Fixes $\alpha$, but the servos
  get $12\times$ less stiff. Gravity sag $\tau_g/K_p$ grows $12\times$, and the lag
  $(K_d/K_p)\dot q = (2\zeta/\omega)\dot q$ goes from 0.18 to 0.64\,s$\cdot\dot q$. Hardware transfer is open.
  \item \textbf{Keep the gains, set $\alpha$ per joint} ($c \approx 3$--$4$ on the arms). The plant is
  untouched and $\alpha$ matches the 1\,Hz run. Tested on the wrist only.
  \item \textbf{Measure actions from another anchor}, which removes the $(q-q_0)$ term. Either the reference,
  $q_{sp} = q_{\text{ref}} + \alpha a$, or the current position, $q_{sp} = q + \alpha a$ (mjlab's
  \texttt{RelativeJointPositionAction}; that one loses the position anchor). Untested; both change the
  paper's action definition.
\end{enumerate}
\textbf{Rule of thumb:} when moving BeyondMimic or SuperMimic to a new robot, check $K_p/\tau_{\max}$
per joint and aim for G1's $\approx 0.5$--$0.7$ per rad ($\alpha \approx 0.35$--$0.55$\,rad). The
frequency alone doesn't tell you this.
\textbf{Limits:} single seed per run; the 1\,Hz runs change plant and scale together; the reward
gradient was not measured directly; simulation only.

\end{document}
